\documentclass[12pt]{extarticle}
\def\activityid{F03-U-v2}\usepackage[letterpaper,margin=.65in,top=.60in,bottom=.65in]{geometry}
\usepackage[T1]{fontenc}\usepackage[default]{sourcesanspro}
\usepackage{microtype,tikz,array,tabularx,fancyhdr,amsmath,amssymb}
\usepackage[hidelinks]{hyperref}\usetikzlibrary{calc,arrows.meta}
\setlength{\parindent}{0pt}\setlength{\parskip}{7pt}\setlength{\headheight}{0pt}\setlength{\footskip}{26pt}
\setlength{\emergencystretch}{2em}\pagestyle{fancy}\fancyhf{}\renewcommand{\headrulewidth}{0pt}\renewcommand{\footrulewidth}{.4pt}
\fancyfoot[L]{\scriptsize Bellingham Math Circle / Week 3 / \activityid}\fancyfoot[R]{\scriptsize \thepage}
\definecolor{ink}{gray}{.12}\definecolor{soft}{gray}{.95}\color{ink}
\newcommand{\PageTitle}[2]{{\fontsize{10}{12}\selectfont\bfseries WEEK 3\quad CODE MACHINES\hfill #1\par}\vspace{3pt}{\fontsize{26}{29}\selectfont\bfseries #2\par}\vspace{2pt}}
\newcommand{\NameLine}{{\small Name: \rule{2.65in}{.4pt}\hfill Date: \rule{1.35in}{.4pt}\par}\vspace{4pt}}
\newcommand{\Task}[2]{\par\vspace{3pt}{\large\bfseries #1}\par #2\par}
\newcommand{\Subhead}[1]{\par\vspace{3pt}{\large\bfseries #1}\par}
\newcommand{\RuleBox}[1]{\par\begingroup\setlength{\fboxsep}{8pt}\colorbox{soft}{\parbox{\dimexpr\linewidth-16pt\relax}{#1}}\endgroup\par}
\newcommand{\WriteBox}[2]{\par\textbf{#1}\par\vspace{2pt}\begin{tikzpicture}\draw[black!40,line width=.5pt] (0,0) rectangle (\dimexpr\linewidth-.5pt\relax,#2);\end{tikzpicture}\par}
\newcommand{\StudentFooter}{\vfill{\small Try it with objects. Explain by pointing, talking, or drawing.}\par}
\newcommand{\LampRing}[3]{\begin{tikzpicture}[x=1in,y=1in]
\foreach \i in {1,...,#1}{\coordinate (v\i) at ({90-360*(\i-1)/#1}:#2);}
\foreach \i in {1,...,#1}{\pgfmathtruncatemacro{\j}{mod(\i,#1)+1}\draw[thick] (v\i)--node[midway,fill=white,inner sep=2pt,font=\small]{\i\j}(v\j);}
\foreach \i in {1,...,#1}{\node[circle,draw,fill=white,minimum size=.52in,inner sep=0pt,font=\bfseries] at(v\i){\i};}
\foreach \i in {#3}{\node[circle,draw,fill=ink,text=white,minimum size=.52in,inner sep=0pt,font=\bfseries] at(v\i){\i};}
\end{tikzpicture}}
\newcommand{\Lights}[1]{\begin{tikzpicture}[x=.55in,y=.55in,baseline=-.10in]\foreach \v [count=\i] in {#1}{\ifnum\v=1\fill (\i,0) circle(.16in);\else\draw (\i,0) circle(.16in);\fi}\end{tikzpicture}}
\newcommand{\ClockRing}[2]{\begin{tikzpicture}[x=1in,y=1in]\draw[black!30] (0,0) circle(#2);\pgfmathtruncatemacro{\n}{#1-1}\foreach\i in {0,...,\n}{\node[circle,draw,fill=white,minimum size=.35in,inner sep=0pt] at({90-360*\i/#1}:#2){\i};}\draw[-{Stealth},thick] (-.35,.15) arc(155:25:.38);\node[font=\small] at(0,-.18){clockwise};\end{tikzpicture}}
\newcommand{\Key}[2]{\begin{center}\renewcommand{\arraystretch}{1.55}\begin{tabular}{c|*{#1}{c}}in&#2\end{tabular}\end{center}}

\begin{document}

\PageTitle{GRADES 4--5 / 1}{How long until every letter returns?}\NameLine
\RuleBox{A substitution key uses A, B, C, D, E once each as inputs and once each as outputs. Apply the key to every letter \textbf{simultaneously}. Repeat the same key each turn.}
\begin{center}\renewcommand{\arraystretch}{1.55}\begin{tabular}{c|ccccc}Input&A&B&C&D&E\\\hline Output&B&C&A&E&D\end{tabular}\end{center}
\Task{1. Test this machine.}{Start with ABCDE. Record successive messages until the \textbf{entire} message first returns. Draw the arrow loops separately.}
\begin{center}\renewcommand{\arraystretch}{1.45}\begin{tabular}{c|p{2in}||c|p{2in}}Turn&Message&Turn&Message\\\hline 0&ABCDE&4&\\1&&5&\\2&&6&\\3&&7&\end{tabular}\end{center}
\WriteBox{The loops, and why their clocks must agree:}{.95in}
\Task{2. Make it slower.}{Can a different five-letter key take longer to return? A five-letter loop is one candidate. Try several keys; keep a record of the loop sizes.}
\WriteBox{My best key, its loops, and its return time:}{1.0in}
\StudentFooter
\newpage
\PageTitle{GRADES 4--5 / 2}{Prove the best is best}
\Task{3. Replace a huge list by a small one.}{Cover the table first. Arrange your five letter cards into loops, working by largest loop: 5, then 4, then 3, then 2, then 1. List how the remaining cards can split without a larger loop. Uncover and compare. A fixed letter is a loop of size 1; rearranging sizes gives no new case.}
\begin{center}\renewcommand{\arraystretch}{1.75}\begin{tabular}{p{2.2in}|p{2.5in}}Loop sizes&First common return time\\\hline 5&\\4 + 1&\\3 + 2&\\3 + 1 + 1&\\2 + 2 + 1&\\2 + 1 + 1 + 1&\\1 + 1 + 1 + 1 + 1&\end{tabular}\end{center}
\Task{4. Explain the classification.}{Why does following arrows eventually bring a letter back to \textbf{itself}, rather than merge into another loop? Where does using each output once matter? Why does your largest-loop list omit no possible split?}
\WriteBox{Why all keys are covered, and the largest return time:}{1.1in}
\textbf{A useful language:} A loop of size 3 returns on turns 3, 6, 9, \ldots. The whole key returns at the first turn that is a multiple of \emph{every} loop size. Mathematicians call this the least common multiple.
\StudentFooter
\newpage
\PageTitle{GRADES 4--5 / 3}{Does the order of two keys matter?}
\RuleBox{Key P swaps A and B and leaves C, D, E alone. Key Q swaps B and C and leaves A, D, E alone. ``P then Q'' means apply P first, then Q to its output.}
\Task{5. Compare two pipelines.}{Start with ABCDE each time. Try P then Q. Reset and try Q then P. Do they agree? One differing letter is enough to disprove agreement.}
\begin{center}\renewcommand{\arraystretch}{1.8}\begin{tabular}{l|p{1.8in}|p{1.8in}}&After first key&After second key\\\hline P then Q&&\\Q then P&&\end{tabular}\end{center}
\WriteBox{An arrow picture explaining the difference:}{.95in}
\Task{6. Repeat the entire pipeline.}{P alone returns in two turns; so does Q. If a turn is now \textbf{P then Q}, how many turns until everything returns? Find its loops.}
\WriteBox{My answer and a loop proof:}{1in}
\textbf{Go further:} Make two different nonidentity keys whose order \emph{does not} matter. What if they move completely separate sets of letters? Can you undo P then Q using the same two keys?
\StudentFooter

\end{document}
